\documentclass[10pt]{article}
% TEMPLATE: code/python/weekly_report.py writes notes/simulation_weekly_optimal_<order>.tex from this
% file (notes/templates/simulation_weekly_optimal.tex), with run_values replaced by run_values_<order> (<order> the sweep order of the run);
% compile that one. Edit the text here.
% the values of one run of the weekly simulation, generated by code/python/weekly_report.py
\input{run_values}
% the four figures, in the folder \RunFigDir of the run (relative to this file)
\newcommand{\figrho}{fig_compare_rho.png}
\newcommand{\figcomp}{fig_weekly_optimal_compartments.png}
\newcommand{\figinf}{fig_weekly_optimal_infectious.png}
\newcommand{\figdeaths}{fig_weekly_optimal_deaths.png}
\usepackage[margin=0.65in]{geometry}
\usepackage{amsmath,amssymb}
\usepackage{graphicx}
\usepackage{xurl}
\usepackage[colorlinks=true,urlcolor=blue]{hyperref}
\pagestyle{empty}
% less space above a paragraph heading
\makeatletter
\renewcommand\paragraph{\@startsection{paragraph}{4}{\z@}{0.8ex \@plus 0.3ex}{-1em}%
  {\normalfont\normalsize\bfseries}}
\makeatother

\begin{document}
\begin{center}
{\large\bfseries The weekly allocation of vaccine: simulation, 2020--2022}\\[2pt]
Authors\footnote{Parts of this note were prepared with the assistance of an AI tool (Claude
Opus 5.5, Anthropic). The authors are solely responsible for all statements.}
\end{center}

\begin{center}
\includegraphics[width=0.36\textwidth]{\RunFigDir\figrho}\hspace{0.04\textwidth}
\includegraphics[width=0.36\textwidth]{\RunFigDir\figcomp}\\[2pt]
\includegraphics[width=0.36\textwidth]{\RunFigDir\figinf}\hspace{0.04\textwidth}
\includegraphics[width=0.36\textwidth]{\RunFigDir\figdeaths}
\end{center}

\paragraph{The simulation.}
The model of Feng, Day and Glasser (J.~Theoret.\ Biol.\ 633, 2026), described in
\texttt{DATA.pdf}, Section~2: 13 compartments for each of $m=51$ locations (50 states and
D.C.) and $n=16$ age groups, $10\,608$ unknowns. Time $t$ is in days from 31 December 2019.
The parameters are the set \texttt{\RunParams} (\texttt{DATA.pdf}, Section~3); the values
that vary in time change at \RunSwitch{}. The system is solved by the Runge--Kutta method of
order 8 (\texttt{DOP853}) from the initial state at $t=\RunTzero$ (\RunTzeroDate): without
vaccination to $t=\RunTvacc$ (\RunTvaccDate), then over the \RunWeeks{} weeks of the dose file to $t=\RunTend$
(\RunTendDate).

Each week the doses of each location are allocated among its age groups so as to minimize
the spectral radius $\rho=\rho(H)$ of the next generation matrix $H$, with at most
$99\%$ of a group vaccinated in a week. The minimization starts from the uniform allocation,
from the rates of the eigenvector strategy, and from the rates of the eigenvector run in
the same week, and keeps the smallest $\rho$; hence $\rho$ never exceeds that of the
eigenvector strategy at the same state. Each sweep of the alternating method minimizes over
\RunOrder{} in turn. The program is
\texttt{code/python/seir7\_ode.py simulate optimal}.

\paragraph{Results.}
With the optimal allocation, $\rho>1$ in \RunRhoOptAbove{} of the \RunWeeks{} weeks, and
some group reaches the bound of $99\%$ in \RunCapWeeks{} weeks. The vaccinated $V$ peak at
$\RunVPeak$ million (\RunVPeakDate), the infectious at $\RunInfPeak$ million
(\RunInfPeakDate) and the hospitalized at $\RunHospPeak$ million (\RunHospPeakDate). The
deaths rise from $\RunDeathsStart$ thousand at $t=\RunTvacc$ to $\RunDeathsEnd$ thousand at
$t=\RunTend$.

% the compartments by age group, one figure per age group (01 = 0-4, ..., 16 = 75+), in a 4 x 4
% grid; they exist only for a run that saved them, otherwise left out
\newcommand{\agefig}[1]{\includegraphics[width=0.24\textwidth]{\RunFigDir fig_weekly_optimal_compartments_age#1.png}}
\IfFileExists{\RunFigDir fig_weekly_optimal_compartments_age01.png}{%
\begin{center}
\agefig{01}\hfill\agefig{02}\hfill\agefig{03}\hfill\agefig{04}\\[2pt]
\agefig{05}\hfill\agefig{06}\hfill\agefig{07}\hfill\agefig{08}\\[2pt]
\agefig{09}\hfill\agefig{10}\hfill\agefig{11}\hfill\agefig{12}\\[2pt]
\agefig{13}\hfill\agefig{14}\hfill\agefig{15}\hfill\agefig{16}
\end{center}}{}

\paragraph{Public data.}
Nine of the 11 inputs come from public sources, each described with its source and script
in \texttt{DATA.pdf}, Section~1: the population $N$ (Census Vintage 2024), the contact
matrix $U$ (Prem, Cook, Jit 2017), the distances and activity factors (Census centers of
population, county estimates, Gazetteer), the contact matrix $C$ (from $U$, $L$, $d$), the
stringency $s_i(t)$ (OxCGRT), the age-specific values and $\beta^0$ (Glasser et al.\ 2023,
Table~2 and Fig.~4), and the days of the variants (CDC variant proportions). Rebuilt from
\texttt{data/Downloaded\_from\_source/} by the scripts named in \texttt{DATA.pdf}, all 11
inputs are identical to those used.

\paragraph{The initial state and the doses.}
The other two inputs are \RunOtherInputs, and the results depend on them:
(1)~\RunInitialText; (2)~\RunDosesText.


\begin{center}
{\large\bfseries Comparison with the eigenvector strategy}
\end{center}

\paragraph{The two strategies.}
The model, the inputs and the weekly doses of each location are those above; the two
strategies differ only in how they allocate these doses among the age groups
(\texttt{DATA.pdf}, Section~2.3):
\begin{itemize}
\item \emph{optimal}: the rates that minimize $\rho$;
\item \emph{eigenvector}: the alternative strategy of Feng, Day and Glasser (2026,
Section~2.2), doses in proportion to the left eigenvector of $H$ at the start of the week,
computed with the rates of the previous week, with the choices (C1)--(C4) described in
\texttt{DATA.pdf}, Section~2.3.
\end{itemize}
Each strategy has its own solution from $t=\RunTvacc$, and $\rho$ is computed each week from
its own state and rates. The two states agree only in the first week, so later the
eigenvector strategy can have the smaller $\rho$. The results are in
\texttt{weekly\_\allowbreak{}simulation/\allowbreak{}results/\allowbreak{}\RunTag/}.

\paragraph{Results of the two strategies.}
\begin{center}
\begin{tabular}{lcc}
 & optimal & eigenvector\\\hline
weeks with $\rho>1$ & \CmpOptimalRhoAbove & \CmpEigenvectorRhoAbove\\
peak of the infectious (millions) & \CmpOptimalInfPeak & \CmpEigenvectorInfPeak\\
peak of the hospitalized (millions) & \CmpOptimalHospPeak & \CmpEigenvectorHospPeak\\
deaths at $t=\RunTend$ (thousands) & \CmpOptimalDeathsEnd & \CmpEigenvectorDeathsEnd
\end{tabular}
\end{center}

% the doses-by-age figures exist only for a run whose two strategies both record the doses by
% age; without them, the figure and the paragraph are left out
\begin{center}
\includegraphics[width=0.48\textwidth]{\RunFigDir fig_compare_infectious.png}\hfill
\includegraphics[width=0.48\textwidth]{\RunFigDir fig_compare_hospitalized.png}\\[4pt]
\includegraphics[width=0.48\textwidth]{\RunFigDir fig_compare_deaths.png}\hfill
\IfFileExists{\RunFigDir fig_compare_doses_age_dose1.png}{%
\includegraphics[width=0.48\textwidth]{\RunFigDir fig_compare_doses_age_dose1.png}}{}
\end{center}

\IfFileExists{\RunFigDir fig_compare_doses_age_dose2.png}{%
\paragraph{Doses by age.} Each dose by age group, summed over the weeks and the locations:
dose 1 above, doses 2 and 3 below.
\begin{center}
\includegraphics[width=0.48\textwidth]{\RunFigDir fig_compare_doses_age_dose2.png}\hfill
\includegraphics[width=0.48\textwidth]{\RunFigDir fig_compare_doses_age_dose3.png}
\end{center}}{}

\end{document}
